\chapter{fast vs strict：研究速度和成交真实性的边界}

\section{本章要解决的困惑}

为什么既有 fast backtest，又有 strict replay？因为它们回答的问题不同。

\begin{longtable}{p{0.20\textwidth}p{0.34\textwidth}p{0.36\textwidth}}
\toprule
模式 & 回答什么 & 不回答什么 \\
\midrule
fast & 因子/参数/entry/cell/capacity 的快速研究 & Bot 响应、真实 order arrival、event log 因果链 \\
strict & 本地交易所里订单如何到达和成交 & 大批量扫因子、快速参数探索 \\
\bottomrule
\end{longtable}

\section{fast 是研究显微镜}

fast 线适合：

\begin{itemize}[leftmargin=2em]
  \item 快速比较 profile；
  \item 检查 spread gate；
  \item 看 daily split；
  \item 看 cell/gamma/capacity；
  \item 研究 fixed60 mid/taker 的差异。
\end{itemize}

fast 线不适合：

\begin{itemize}[leftmargin=2em]
  \item 验证 TCP bridge；
  \item 验证 latency；
  \item 验证 event log；
  \item 验证 Bot 的实时处理是否等价。
\end{itemize}

\section{strict 是本地交易所}

strict replay 适合：

\begin{itemize}[leftmargin=2em]
  \item 检查 order intent 到 fill 的因果链；
  \item 验证 observed quote 和 arrival quote；
  \item 检查 fill price 是否符合模型；
  \item 记录 portfolio state；
  \item 让 Monitor 读真实 run artifact。
\end{itemize}

strict 不是为了扫 1000 个参数。那会慢，而且容易把工程验证当成研究优化。

\section{两个模式怎样对齐}

最小对齐问题是：

\begin{lstlisting}
same market date
same policy profile
same entry clock
same admission threshold
same capacity profile
same fill approximation
\end{lstlisting}

如果 fast 用 \code{fixed60_mid}，strict 用 top-of-book taker fill，两者就不能直接比较净收益。至少要说明 execution approximation。

\section{什么时候必须 strict}

以下情况必须用 strict 或 audit：

\begin{itemize}[leftmargin=2em]
  \item 改了 Bot 事件处理；
  \item 改了 Runner fill/latency/order code；
  \item 改了 event contract；
  \item 要证明 Monitor 显示的是事实；
  \item 研究从 top-of-book taker 切到 L2-depth 或 maker。
\end{itemize}

\section{什么时候 fast 就够}

以下情况先用 fast：

\begin{itemize}[leftmargin=2em]
  \item 只是看某个 entry gate 有没有基本信号；
  \item 比较 q70/q80 admission；
  \item 比较 gamma profile；
  \item 检查 cell 分布；
  \item 生成下一步研究假设。
\end{itemize}

\section{常见误解}

\begin{itemize}[leftmargin=2em]
  \item \term{误解一：strict 一定更真实，所以所有研究都 strict。}不对。strict 慢，适合验收，不适合大范围探索。
  \item \term{误解二：fast 赚了就能上线。}不对。fast 只是近似，必须看 taker executable、daily split 和 strict audit。
  \item \term{误解三：两者收益差就是 bug。}不一定。先检查 profile、clock、fill model 是否一致。
\end{itemize}

\checkpoint{fast 负责快，strict 负责真。把两者混成一个东西，代码就会越来越难读。}

